\documentclass[11pt]{article}
\usepackage{amsmath,amssymb,mathtools}
\usepackage{geometry}
\usepackage{microtype}
\geometry{margin=1in}

\title{Technical Report: Machine Learning Classification and Regression Pipeline}
\author{}
\date{}

\begin{document}
\maketitle

\begin{abstract}
This report provides a technical analysis of a Python implementation that applies classical
machine-learning methods to student performance data. The program implements linear
regression, several binary classification models, logistic regression decision-boundary
analysis, and a manual matrix-based classifier formulation. The analysis covers the
computational workflow, mathematical foundations, implementation details, and
reproducibility considerations.
\end{abstract}

\section{Introduction}
The project implements a classical machine-learning workflow in Python using Pandas,
NumPy, Matplotlib, Seaborn, and Scikit-learn. The input is an Excel dataset containing
student-related features. The implementation models both continuous outcomes and binary
pass/fail decisions using linear regression and classification algorithms.

\section{Project Overview}
The computational pipeline consists of:
\begin{enumerate}
\item Loading the Excel input file.
\item Preparing feature matrices by removing non-predictive columns.
\item Training a linear regression model for continuous prediction.
\item Training Ridge, SGD, and Logistic Regression classifiers.
\item Computing evaluation measures such as accuracy, precision, and confusion matrices.
\item Analysing a two-feature logistic regression decision boundary.
\item Reconstructing a classifier using explicit matrix operations.
\item Packaging generated artifacts for reproducibility.
\end{enumerate}

\section{Data Processing}
The implementation reads the Excel file into a Pandas dataframe. Columns representing
the target value and metadata are removed before model fitting. For classification, the
categorical pass label is converted to a binary representation:
\[
y_i =
\begin{cases}
0,&\text{if the student fails},\\
1,&\text{if the student passes}.
\end{cases}
\]

\section{Linear Regression}
The regression component learns a linear relationship:
\[
\hat{y}=w^Tx+b,
\]
where $x$ is the feature vector, $w$ contains the learned coefficients, and $b$ is the
intercept. The implementation uses Scikit-learn's LinearRegression estimator and exports
the learned parameters.

\section{Classification Models}
Three classifiers are trained.

\subsection{Ridge Classifier}
Ridge classification adds regularization to a linear decision function. A typical
regularized objective can be written as:
\[
\min_w ||Xw-y||^2_2+\lambda||w||^2_2.
\]

\subsection{Logistic Regression}
Logistic regression models the probability of the positive class as:
\[
P(y=1|x)=\sigma(w^Tx+b),
\]
where:
\[
\sigma(z)=\frac{1}{1+e^{-z}}.
\]
The decision boundary occurs when the model score is zero:
\[
w^Tx+b=0.
\]

\subsection{SGD Classifier}
The stochastic gradient descent classifier updates its parameters iteratively using
sampled training examples. Its optimization behavior depends on the Scikit-learn
configuration used by the estimator.

\section{Evaluation Procedure}
The implementation calculates accuracy:
\[
\mathrm{Accuracy}=\frac{TP+TN}{TP+TN+FP+FN},
\]
precision:
\[
\mathrm{Precision}=\frac{TP}{TP+FP},
\]
and confusion matrices for the classifiers.

\section{Decision Boundary Analysis}
A logistic regression model is trained using only midterm and final exam scores. The
decision boundary is derived from:
\[
b+w_1x_1+w_2x_2=0.
\]
Therefore:
\[
x_2=-\frac{b}{w_2}-\frac{w_1}{w_2}x_1 .
\]

\section{Manual Matrix Formulation}
The project also demonstrates a manual linear classifier calculation. A bias term is
inserted into the feature matrix:
\[
X=[1,x_1,x_2].
\]
The implementation computes matrix products involving:
\[
(X^TX)^{-1}X^T
\]
and derives classifier weights through explicit numerical operations.

\section{Implementation Details}
The program is adapted for Google Colab execution. It automatically installs dependencies,
handles the Excel input file, creates output directories, stores numerical results, exports
tables, and creates a compressed archive containing the generated artifacts.

\section{Reproducibility Considerations}
The implementation supports reproducibility through explicit input handling, consistent
file organization, and export of learned parameters. However, no fixed random seeds or
train/test separation are defined in the source code; therefore, the reported metrics
reflect evaluation on the fitted dataset rather than performance on an independent test set.

\section{Limitations}
The implementation is educational and demonstrates several algorithmic concepts. Because
the models are evaluated on the same data used for training, the resulting metrics should
not be interpreted as estimates of generalization performance. Additional validation
procedures would be required for scientific benchmarking.

\section{Conclusion}
This project demonstrates a coherent classical machine-learning workflow that combines
regression, classification, visualization, mathematical reconstruction, and reproducible
artifact generation. It provides a practical connection between mathematical formulations
and executable machine-learning procedures.

\end{document}
